% Journal-neutral revised manuscript after MEAS-D-26-10542 review.
\documentclass[12pt]{article}

\usepackage[utf8]{inputenc}
\usepackage[T1]{fontenc}
\usepackage{lmodern}
\usepackage{amsmath,amssymb,amsthm}
\usepackage{booktabs}
\usepackage{array}
\usepackage{graphicx}
\usepackage{algorithm}
\usepackage{algorithmic}
\usepackage{enumitem}
\usepackage[margin=1in]{geometry}
\usepackage{lineno}
\usepackage[breaklinks=true,hidelinks]{hyperref}

\newtheorem{proposition}{Proposition}
\renewenvironment{proof}[1][Proof]{\noindent\textit{#1.}\space}{\hfill$\blacksquare$\par}

\begin{document}
\title{A Causal Evaluation of PCA-Guided Sensor-Rate Allocation under Hard Communication Budgets}
\author{Anonymous Author(s)\\Anonymized for double-blind review}
\date{}
\maketitle

\begin{abstract}
Bandwidth-constrained multichannel measurement requires deciding how frequently
each channel should be transmitted before the corresponding measurements are
available. We study a causal rate-allocation rule that converts the diagonal of
an incremental rank-$k$ PCA covariance approximation into per-channel rates.
Rates for each window are computed only from preceding reconstructed windows;
the first window uses uniform allocation, missing values are reconstructed by
causal zero-order hold, and integer sampling quotas enforce the communication
budget in every window. The allocation rule requires no task labels and makes
the information pattern, reconstruction delay, and realized communication cost
explicit. We evaluate it using independent simulation runs or official
chronological train/test partitions and report measurement distortion,
point-wise and event-level detection, false alarms, detection delay, and actual
bandwidth. At 50\% bandwidth, PCA-guided allocation obtains TEP weighted F1 0.675, compared with 0.707 for Uniform and 0.728 for full data. SMD-1-1 and PSM additionally expose dependence on the downstream detector and threshold. No universal superiority is claimed.

The results show why causal timing, measurement distortion, and downstream
detection must be evaluated separately.
\end{abstract}

\noindent\textbf{Keywords:} adaptive sampling; sensor-rate allocation;
incremental PCA; communication constraints; causal reconstruction; edge measurement

\linenumbers

\section{Introduction}

Distributed measurement systems often collect more channel samples than a
sensor link or edge uplink can transmit. Uniform downsampling is simple, but it
does not account for differences among channels. A rate allocator can instead
assign a rate $r_j$ to each of $d$ channels while respecting a total budget
$B$. The decision is operationally meaningful only if it is made before the
samples governed by that decision are observed.

Prior work includes event-triggered temporal sampling
\cite{miskowicz2006send,gupta2018sbl}, correlation-aware sensor
selection \cite{bacciu2016unsupervised,ghosh2021adaptive}, and online subspace
tracking \cite{weng2003ccipca,balzano2010grouse,yang2018historypca}. These
areas address related questions, but binary channel selection, uniform temporal
sampling, and proportional per-channel rate allocation are not interchangeable.
The term \emph{sensor triage}, used in an earlier version of this work, means
only proportional sensor-rate allocation here.

This study asks whether the channel contributions to an online low-rank
covariance approximation provide a useful unsupervised signal for allocating a
hard communication budget. PCA itself is not presented as novel. The specific
contribution is a causal mapping from incremental PCA covariance contributions
to next-window integer sampling quotas, together with an evaluation protocol
that exposes its measurement and detection trade-offs.

The contributions are:
\begin{enumerate}[leftmargin=*]
\item a causal system and timing model in which rates for window $t$ depend only
on information available through window $t-1$;
\item a dimensionally consistent score equal to the diagonal of the rank-$k$
PCA covariance approximation;
\item a largest-remainder quota rule that enforces the aggregate budget in each
window rather than only in expectation; and
\item an evaluation based on independent runs or official chronological splits,
with distortion, communication, point-level, event-level, and delay metrics.
\end{enumerate}

\section{Related work}

\begin{table}[t]
\centering
\small
\caption{Relationship to the closest methodological families. ``Rate'' means
proportional per-channel rates rather than binary inclusion.}
\resizebox{\textwidth}{!}{%
\begin{tabular}{@{}p{0.16\textwidth}p{0.14\textwidth}p{0.16\textwidth}p{0.18\textwidth}p{0.25\textwidth}@{}}
\toprule
Approach & Decision & Main signal & Strength & Difference from this study \\
\midrule
Send-on-delta \cite{miskowicz2006send} & transmit on change & local deviation & event responsiveness & observes current value before deciding \\
Bacciu \cite{bacciu2016unsupervised} & select feature subset & cross-correlation & redundancy filtering & does not assign temporal quotas \\
Ghosh et al. \cite{ghosh2021adaptive} & active sensor set & UCB, correlation, energy & predicts inactive sensors & energy-aware selection with GPR \\
Gupta and De \cite{gupta2018sbl} & active set and cycle rate & PCA-SBL recovery & feedback on sensing error & selection/recovery optimization \\
CCIPCA \cite{weng2003ccipca} & update subspace & covariance structure & incremental tracking & supplies no allocation rule \\
Present rule & next-window quotas & covariance diagonal & explicit sample cap & descriptive unsupervised proxy \\
\bottomrule
\end{tabular}%
}
\label{tab:positioning}
\end{table}

Send-on-delta triggers transmission when the current observation differs
sufficiently from the previous transmitted value \cite{miskowicz2006send}.
This information pattern is natural for a locally observed sensor but differs
from deciding an acquisition schedule before values arrive. Bacciu's
cross-correlation filter selects nonredundant time-series features
\cite{bacciu2016unsupervised}. Ghosh et al. select an active sensor set using
upper-confidence-bound learning and energy information, and reconstruct
inactive parameters with Gaussian process regression \cite{ghosh2021adaptive}.
These methods address relevance or redundancy more directly than a marginal
covariance-diagonal score.

A particularly close predecessor is the PCA-SBL framework of Gupta and De
\cite{gupta2018sbl}. It combines PCA-based sparse signal recovery, energy-aware
sensor selection, and feedback that predicts the next measurement cycle's
sampling requirement. Consequently, neither combining PCA with adaptive sensing
nor making a next-cycle decision is novel by itself. The present study instead
isolates a simpler proportional covariance-participation rule under an explicit
integer sample cap and common causal reconstruction. It does not establish
superiority over these complete selection and recovery systems.

Recent sensor-sampling work extends this landscape. CLARA combines node
correlation, clustering, sampling-rate adaptation, and fault tolerance in
periodic wireless sensor networks \cite{harb2024clara}. Its network-oriented
scope motivates correlation-aware allocation, but our replay experiments do
not model network clustering or node failures. CoSS jointly selects sensors
and sampling rates using weights learned within a human-activity recognition
network \cite{liu2026coss}. This is a task-trained selection mechanism; the
present allocator uses no class labels to choose quotas. Huang et al. formulate
multi-sensor sampling as a Markov decision process and train a deep Q-network
using a reward that combines information, energy cost, and redundancy
\cite{huang2025dqn}. Their formulation highlights the importance of specifying
both the information available at decision time and the resource objective.
These recent methods are relevant alternatives, not methods against which the
present results establish an empirical ranking.

Online subspace methods such as Oja's rule, CCIPCA, and GROUSE can replace the
incremental SVD update \cite{oja1982simplified,weng2003ccipca,balzano2010grouse}.
They change subspace estimation and computational cost, but do not by themselves
specify a communication policy. The method studied here supplies no new
process-monitoring statistic or control limit; downstream detectors are
measurement consumers used to evaluate the allocation rule.

\section{Method}

\subsection{System, budget, and information pattern}

Let $X_t\in\mathbb{R}^{w\times d}$ denote the full but initially unobserved
measurements in window $t$. Before $X_t$ arrives, a controller announces rates
$r_{t,j}\in[r_{\min},1]$. The sensing or transmission mechanism retains a
subset $\Omega_t$ and the receiver observes only
$P_{\Omega_t}(X_t)$. The aggregate hard constraint is
\begin{equation}
 |\Omega_t|\leq \lfloor Bwd\rfloor .
\label{eq:hardbudget}
\end{equation}
The controller reconstructs $\widehat X_t$ and then updates the PCA state.
Thus $r_t$ is measurable with respect to reconstructed observations through
$t-1$ and cannot depend on unseen values in $X_t$.

The first window is a uniform bootstrap, $r_{1,j}=B$. This cost is included in
all reported communication measurements. The strict acquisition-constrained
model updates PCA with $\widehat X_t$. A separate uplink-constrained variant,
in which a local gateway legitimately observes full-rate $X_t$, may update from
$X_t$ but must be reported as a different information model.

Windows are non-overlapping. This prevents one raw observation from affecting
multiple nominal decisions and makes payload accounting direct. A final short
window receives its own integer quota. No model, reconstruction state, or
window crosses an independent machine or simulation-run boundary.

\subsection{PCA covariance-contribution score}

After processing window $t$, incremental PCA gives eigenpairs
$(\lambda_i,v_i)$ for $i=1,\ldots,k$. We define
\begin{equation}
 s_{t,j}=\sum_{i=1}^{k}\lambda_i v_{ij}^{2}.
\label{eq:score}
\end{equation}
This is the $j$th diagonal element of the rank-$k$ approximation
\begin{equation}
 C_k=\sum_{i=1}^k\lambda_i v_iv_i^\top.
\end{equation}
It measures how much retained
subspace variance is expressed through channel $j$. It is not a supervised
measure of fault relevance and does not select one representative from a
correlated pair. This limitation follows from the score definition.

The implemented method first normalizes the raw contributions:
\begin{equation}
 u_{t,j}=s_{t,j}/\sum_l s_{t,l}.
\end{equation}
If the total retained variance is zero or nonfinite, $u_t$ is uniform. This
fallback prevents a constant bootstrap from poisoning subsequent updates.
For $0\leq\rho<1$, smoothing operates on normalized participation:
\begin{equation}
 \bar u_t=\rho\bar u_{t-1}+(1-\rho)u_t,
 \qquad \bar u_1=u_1.
\end{equation}
The special value $\rho=1$ denotes the cumulative mean of the normalized score
vectors. Normalization before smoothing prevents a high-energy window from
receiving additional weight solely through its covariance trace; this differs
from smoothing raw covariance contributions and then normalizing.

\subsection{Rates, integer quotas, and acquisition}

Set $p_{t,j}=\bar u_{t,j}/\sum_l\bar u_{t,l}$. The update after window $t$
sets the next window's target rates to
\begin{equation}
 r_{t+1,j}=r_{\min}+p_{t,j}(Bd-r_{\min}d).
\label{eq:rates}
\end{equation}
After clipping, excess is repeatedly redistributed among unsaturated channels
until the feasible target budget is exhausted. We floor
the target counts $wr_{t+1,j}$, then use largest remainders to reach exactly
$\lfloor Bwd\rfloor$ samples. Positions are drawn without replacement. The minimum rate is a continuous
target floor: integer rounding can reduce an individual realized rate below
that floor, especially in a short final window. The aggregate cap is exact.

The receiver uses zero-order hold. Leading missing values use the last value
from the preceding window; at cold start they use zero after train-fitted
standardization. Linear interpolation is excluded from the primary online
protocol because it requires future samples. The software retains it for explicitly offline use, but no such analysis is
included in the primary or supplementary evidence here.

\begin{algorithm}[t]
\caption{Causal PCA-guided rate allocation}
\label{alg:causal}
\begin{algorithmic}[1]
\STATE Set $r_{1,j}\gets B$ for all $j$; initialize PCA
\FOR{$t=1,2,\ldots$}
  \STATE Announce $r_t$ before observing window $X_t$
  \STATE Convert $r_t$ to integer quotas summing to $\lfloor Bwd\rfloor$
  \STATE Acquire samples $P_{\Omega_t}(X_t)$ under those quotas
  \STATE Causally reconstruct $\widehat X_t$ by zero-order hold
  \STATE Update incremental PCA using $\widehat X_t$
  \STATE Compute and smooth scores using Eq.~\eqref{eq:score}
  \STATE Allocate $r_{t+1}$ using Eq.~\eqref{eq:rates}
\ENDFOR
\end{algorithmic}
\end{algorithm}

\subsection{Complexity}

For a window of length $w$, $d$ channels, and $k$ retained components, the
incremental SVD update dominates computation. No controller-specific timing or microcontroller-readiness claim is made
from operation counts alone. End-to-end claims require timing the acquisition, reconstruction,
update, allocation, and communication-control path on named hardware.

\section{Properties and scope of the score}

\begin{proposition}[Covariance-diagonal identity]
For a positive semidefinite covariance matrix
$C=\sum_{i=1}^{d}\lambda_i v_iv_i^\top$, the score in
Eq.~\eqref{eq:score} satisfies
\begin{align}
 s_j&\geq0, & \sum_{j=1}^{d}s_j&=\sum_{i=1}^{k}\lambda_i,\\
 C_{jj}-s_j&=\sum_{i=k+1}^{d}\lambda_i v_{ij}^{2}\geq0.
\end{align}
For $k=d$, $s_j=C_{jj}$ and the raw PCA score equals the diagonal variance of the same covariance matrix.
\end{proposition}
\begin{proof}
Non-negativity follows from $\lambda_i\geq0$ and squared loadings. Summing over
$j$ and using $\sum_jv_{ij}^2=1$ gives the trace identity. Subtracting the
truncated eigendecomposition from the full diagonal gives the residual. For
$k=d$ the residual is zero.
\end{proof}

This proposition supplies a rationale for the score but not an optimality
claim. In particular, two exchangeable correlated channels receive equal
scores. The previous claim that squared PCA loadings necessarily concentrate
bandwidth on one member of a correlated pair is false and has been removed.

\begin{proposition}
Assume $0\leq r_{\min}\leq B$. Largest-remainder integerization of feasible
target rates produces nonnegative quotas $q_{t,j}\leq w$ satisfying
$\sum_jq_{t,j}=\lfloor Bwd\rfloor$ whenever that many samples are available.
\end{proposition}
\begin{proof}
Flooring produces an integer total no larger than the target. Each remaining
unit is assigned only to a channel below its capacity. The procedure stops at
the prescribed integer total, which is at most $wd$ because $B\leq1$.
\end{proof}

If incremental PCA converges under a stationary distribution and its retained
trace has a positive limit, the normalized score converges by continuity.
For $\rho<1$, its exponential average converges to the same limit; for
$\rho=1$, the cumulative mean converges by Ces\`aro convergence.
No convergence claim is made for abrupt nonstationarity without additional
assumptions.

\section{Experimental protocol}

\subsection{Datasets and splits}

All 20 TEP fault types are evaluated. The 52 transmitted process channels
include 41 measured and 11 manipulated variables; the budget counts values
rather than assuming identical physical sensing costs. Simulations are split by simulation run;
no window crosses a run, and online state is reset between runs
\cite{downs1993tep,rieth2017tep}. SMD
and PSM use their official normal-training and labeled-test streams
\cite{su2019omnianomaly,abdulaal2021ransyn}. Scaling is
fit on training observations only. Random stratification of time points is not
used. Hyperparameters are global unless selected exclusively from a validation
subset on the training side.

\begin{table}[t]
\centering
\caption{Corrected dataset protocol. TEP counts are per fault type; SMD uses
machine-1-1 and is not presented as the entire SMD collection.}
\begin{tabular}{@{}lrrrl@{}}
\toprule
Dataset & Channels & Train & Test & Boundary \\
\midrule
TEP & 52 & $5\times500$ & $5\times500$ & simulation run \\
SMD-1-1 & 38 & 28,479 & 28,479 & official stream \\
PSM & 25 & 132,481 & 87,841 & official stream \\
\bottomrule
\end{tabular}
\label{tab:datasets}
\end{table}

TEP faulty trajectories contain a normal prefix followed by an introduced
disturbance. Samples before onset are labeled normal instead of assigning the
fault label to the entire trajectory. Both partitions use the released Training RData files: runs 1--5 train the
classifier and runs 21--25 provide final evaluation for every fault type.
The label convention is normal for sample indices below 20 and faulty from
index 20 onward; it is recorded explicitly in the experiment manifest. SMD machine-1-1 and
PSM retain their released order; normal training streams are not mixed with
labeled test events. Using one SMD machine is explicit because concatenating
machines would create artificial boundaries in a supposed single stream.

The aggregate MSL arrays used during an earlier revision are excluded from
primary evidence. Their provenance does not establish the original entity
boundaries; the source contains per-stream telemetry and encoded command inputs
rather than 55 interchangeable physical sensors \cite{telemanomdata}.
Treating the aggregate array as one continuous sensor stream would conflate
independent histories and an inadequately specified communication model.

Allocation, PCA, and reconstruction state start fresh at every official train
or test partition and every independent TEP run. The test stream therefore
starts with a counted uniform bootstrap, rather than retaining part of the
training controller's state. Scaling and trained downstream detectors are
retained from training. Full-rate training observations are available offline
for the common scaler and fixed-detector reference. They are calibration inputs,
not observations available to the online test allocator. Reported payload totals
cover test acquisition; separate training ledgers are supplied. The budget
therefore characterizes deployment replay, not total lifecycle acquisition cost. Partial final windows use the latest available rates;
only complete windows update PCA.

Nonfinite numeric values are replaced before scaling. A StandardScaler is fit
only to training observations and then applied without refitting to test data.
The standardized value zero is consequently a train-derived mean and supplies
the cold-start value when a channel has not delivered an observation. Windows
do not overlap, labels are unavailable to the allocator, and no parameter is
selected from final test labels.

\subsection{Comparators and downstream protocols}

The primary comparators are Uniform allocation and two prior-window rules:
Variance and change Threshold. All receive the same information, integer
budget, bootstrap, and causal reconstruction. The primary study contains only these constrained rules and Full Data.
Send-on-delta and supervised feature selection are discussed as related work,
not represented as evaluated peers.

For labeled TEP runs, a random forest is evaluated in two ways: retrained for
each reconstructed training set, and fixed after training on full-rate data
\cite{breiman2001random}. Official normal-only anomaly-training datasets use an
Isolation Forest with a threshold determined from training scores. The same
retrained-versus-fixed distinction is reported. These protocols separate
end-to-end adaptation from distortion of an existing detector.

Variance allocation computes per-channel variance from the preceding causally
reconstructed window. Threshold allocation uses preceding-window mean absolute
changes and prioritizes channels above the median change. Uniform allocation
assigns the same target rate to every channel. All three use the same integer
quota routine and zero-order hold as the PCA rule. Full Data is a reference,
not a constrained competitor. Budget realization and reconstruction are held constant. However, PCA uses
accumulated subspace information and smoothing, whereas Variance and Threshold
use one preceding window. The comparison therefore evaluates complete allocation
rules rather than isolating a PCA effect independently of temporal memory.

\subsection{Scope of the comparisons with recent methods}
The experimental question is whether the covariance-participation rule improves
on elementary allocation rules when each method receives only previously
reconstructed samples and must obey the same next-window integer sample cap.
Uniform, Variance, and Threshold are transparent reference rules, not a claim
to cover the state of the art. Full Data measures the cost of restricted
acquisition and is not a budget-matched competitor.

We do not report numerical comparisons with CLARA, CoSS, or the DQN sampling
method. The reasons, and the resulting limits on interpretation, are as follows.
\begin{enumerate}[leftmargin=*]
\item CLARA's clustering and fault-tolerance setting \cite{harb2024clara} would
require an explicit mapping from network nodes and their relationships to
benchmark channels, together with a network/failure model. Our benchmark
protocol supplies neither. Removing these components would evaluate an
adaptation, whose results could not be attributed to the complete method.
\item CoSS \cite{liu2026coss} learns sensor and rate rankings with a downstream
classification network. A faithful TEP comparison would require a separate
task-trained model, validation-only selection, and accounting for the data
used during selection. The normal-only SMD and PSM training splits do not
supply the labeled activity-classification task used by CoSS. Such extensions
are possible, but are not experiments performed in this study.
\item The DQN method \cite{huang2025dqn} uses a state containing observations,
remaining power, and sampling history, with a reward trading information
against energy and redundancy. The present replay protocol has no calibrated
power model and counts acquired scalar values rather than a reward penalty.
A comparison would need to specify how states are formed from acquired values,
charge all measurements used in learning and decisions, and constrain actions
to the same per-window cap. We have not implemented or validated that adaptation.
\end{enumerate}
These are scope-based reasons for the omitted experiments, not evidence that
recent methods cannot be compared or would perform worse. They leave the
empirical contribution limited to the implemented reference rules. In particular,
TEP labels could support a separate supervised selection study, and clustered
or learned policies could be assessed under an explicitly extended protocol.
A broader comparative study remains necessary before claiming competitive
performance against contemporary sensor-selection or learned-sampling systems.

\subsection{Metrics and repetition}

We report payload values and bytes, realized bandwidth, per-window budget
excess, MAE, RMSE, normalized per-channel RMSE, precision, recall, F1, false
positive rate, specificity, event recall, false-alarm events per 1000 samples,
and detection delay. Point adjustment is not used. Experiments are repeated
with five fixed joint seeds that control acquisition masks, detector randomness,
and (for Isolation Forest) the fitting subset. Means and standard deviations
therefore describe joint experimental variation, not acquisition alone or
population uncertainty. Independent TEP simulation
runs provide substantive replication; the official anomaly benchmarks are
single-stream case studies, so no population-level inferential superiority is
claimed for them.

For original and reconstructed streams $X$ and $\widehat X$, distortion is
summarized by
\begin{equation}
 \mathrm{RMSE}=\sqrt{\frac{1}{nd}\sum_{t=1}^{n}\sum_{j=1}^{d}
 (X_{tj}-\widehat X_{tj})^2}.
\end{equation}
TEP RMSE pools squared errors over all test runs before taking the square root.
Per-channel RMSE is normalized by the pooled test range and averaged; constant
channels use a denominator of one. Test ranges are used only for retrospective
error reporting, never for allocation, training, or threshold selection.
An event is a maximal contiguous positive-label interval. Event recall counts
an event when at least one unadjusted positive prediction falls inside it.
Detection delay is measured from event start to the first such prediction. A
predicted positive interval overlapping no true event is a false-alarm event.
Each independent TEP run is scored separately for events before pooling counts
and delay samples. Delay is conditional on detection and expressed in samples;
missed-event counts are reported alongside it. If no event is detected, delay
is undefined. Per-1000-sample alarm rates do not imply a common physical time
scale across datasets. These definitions avoid retrospective point adjustment.

The global allocation configuration is $k=10$, $w=50$, $\rho=0.95$, and
$r_{\min}=0.05$. Anomaly benchmarks use $B\in\{0.3,0.5,0.7\}$; the larger
all-20-fault TEP experiment uses the primary comparison $B=0.5$.
Sampling seeds are 42, 123, 456,
789, and 1024. TEP uses a 50-tree class-balanced random forest. Normal-only
benchmarks use a 100-tree Isolation Forest fitted to at most 20,000 normal
training observations; its threshold is the 99th percentile of training
anomaly scores. These values are fixed across datasets and are not selected
from test outcomes.

Every acquisition call records the actual retained-mask count, available values,
allowed integer quota, and excess. Results sum those ledgers, and a verifier
independently reconciles every test total. Payload bytes assume four bytes per
transmitted numeric value; they do not measure a network encoding or packet
traffic. Full Data has a separately counted full observation mask.

Threshold sensitivity evaluates the predeclared training-score quantiles
0.95, 0.975, 0.99, 0.995, and 0.999 for every anomaly method and seed. The primary
quantile remains 0.99; no test outcome chooses a preferred threshold. Complete
point, event, distortion, and communication results for both detector protocols
and all budgets appear in the supplementary material.

\section{Results}

\IfFileExists{tables/causal_summary.tex}{
The causal, mask-audited results depend on dataset, budget, and metric. All constrained methods meet their integer sample cap in every recorded window. Payload is an assumed four-byte representation, not measured network traffic.

\begin{table*}[t]
\centering
\caption{Corrected causal results at 50\% bandwidth. Values are mean $\pm$ standard deviation across five joint experimental seeds. Full Data is shown as a detector reference. No point adjustment is used.}
\resizebox{\textwidth}{!}{%
\begin{tabular}{@{}llcccc@{}}
\toprule
Dataset & Method & Retrained F1 & Fixed F1 & Event recall & RMSE \\
\midrule
SMD-1-1 & PCA-Triage & 0.390 $\pm$ 0.017 & 0.401 $\pm$ 0.012 & 0.700 $\pm$ 0.168 & 9.024 $\pm$ 2.230 \\
 & Variance & 0.407 $\pm$ 0.019 & 0.406 $\pm$ 0.011 & 0.825 $\pm$ 0.143 & 1.258 $\pm$ 0.246 \\
 & Threshold & 0.393 $\pm$ 0.013 & 0.396 $\pm$ 0.013 & 0.725 $\pm$ 0.163 & 13.987 $\pm$ 17.107 \\
 & Uniform & 0.395 $\pm$ 0.008 & 0.400 $\pm$ 0.010 & 0.850 $\pm$ 0.137 & 83.863 $\pm$ 19.805 \\
 & Full Data & 0.409 $\pm$ 0.010 & 0.409 $\pm$ 0.010 & 0.875 $\pm$ 0.125 & 0.000 $\pm$ 0.000 \\
\addlinespace
PSM & PCA-Triage & 0.031 $\pm$ 0.001 & 0.035 $\pm$ 0.007 & 0.139 $\pm$ 0.026 & 0.222 $\pm$ 0.011 \\
 & Variance & 0.032 $\pm$ 0.001 & 0.035 $\pm$ 0.007 & 0.153 $\pm$ 0.022 & 0.193 $\pm$ 0.011 \\
 & Threshold & 0.032 $\pm$ 0.002 & 0.034 $\pm$ 0.007 & 0.125 $\pm$ 0.028 & 0.228 $\pm$ 0.014 \\
 & Uniform & 0.033 $\pm$ 0.004 & 0.035 $\pm$ 0.007 & 0.147 $\pm$ 0.032 & 0.237 $\pm$ 0.018 \\
 & Full Data & 0.036 $\pm$ 0.007 & 0.036 $\pm$ 0.007 & 0.214 $\pm$ 0.016 & 0.000 $\pm$ 0.000 \\
\bottomrule
\end{tabular}%
}
\label{tab:causal-anomaly}
\end{table*}

\begin{table}[t]
\centering
\caption{Point-F1 sensitivity to the hard communication budget. The comparator column is the strongest non-PCA constrained method at that budget.}
\begin{tabular}{@{}lrrlr@{}}
\toprule
Dataset & Budget & PCA-Triage & Best comparator & F1 \\
\midrule
SMD & 0.3 & 0.396 & Variance & 0.406 \\
SMD & 0.5 & 0.390 & Variance & 0.407 \\
SMD & 0.7 & 0.396 & Threshold & 0.409 \\
\addlinespace
PSM & 0.3 & 0.030 & Threshold & 0.033 \\
PSM & 0.5 & 0.031 & Uniform & 0.033 \\
PSM & 0.7 & 0.032 & Threshold & 0.036 \\
\bottomrule
\end{tabular}
\label{tab:budget-sensitivity}
\end{table}

\begin{table*}[t]
\centering
\caption{Run-separated TEP results at 50\% bandwidth, with the fault onset labeled within each run. Values are mean $\pm$ standard deviation across five joint experimental seeds.}
\resizebox{\textwidth}{!}{%
\begin{tabular}{@{}lccccc@{}}
\toprule
Method & Retrained weighted F1 & Fixed weighted F1 & Macro F1 & Binary F1 & RMSE \\
\midrule
PCA-Triage & 0.675 $\pm$ 0.003 & 0.671 $\pm$ 0.004 & 0.684 $\pm$ 0.003 & 0.932 $\pm$ 0.002 & 0.270 $\pm$ 0.002 \\
Variance & 0.677 $\pm$ 0.002 & 0.673 $\pm$ 0.004 & 0.686 $\pm$ 0.003 & 0.931 $\pm$ 0.002 & 0.261 $\pm$ 0.002 \\
Threshold & 0.643 $\pm$ 0.005 & 0.644 $\pm$ 0.004 & 0.651 $\pm$ 0.005 & 0.924 $\pm$ 0.003 & 0.383 $\pm$ 0.002 \\
Uniform & 0.707 $\pm$ 0.003 & 0.686 $\pm$ 0.002 & 0.716 $\pm$ 0.003 & 0.937 $\pm$ 0.001 & 0.483 $\pm$ 0.001 \\
Full Data & 0.728 $\pm$ 0.001 & 0.728 $\pm$ 0.001 & 0.740 $\pm$ 0.002 & 0.958 $\pm$ 0.001 & 0.000 $\pm$ 0.000 \\
\bottomrule
\end{tabular}%
}
\label{tab:causal-tep}
\end{table*}

\begin{table*}[t]
\centering
\caption{Exact fault-type recall for all 20 TEP disturbances at 50\% bandwidth, averaged across five joint experimental seeds. Pre-onset samples are labeled normal; confusing one fault for another is counted as incorrect here.}
\resizebox{0.8\textwidth}{!}{%
\begin{tabular}{@{}rrrrrr@{}}
\toprule
Fault & PCA-Triage & Variance & Threshold & Uniform & Full Data \\
\midrule
1 & 0.981 & 0.980 & 0.983 & 0.980 & 0.986 \\
2 & 0.969 & 0.970 & 0.970 & 0.970 & 0.973 \\
3 & 0.158 & 0.155 & 0.113 & 0.167 & 0.145 \\
4 & 0.958 & 0.945 & 0.952 & 0.975 & 0.977 \\
5 & 0.740 & 0.719 & 0.624 & 0.765 & 0.918 \\
6 & 0.997 & 0.997 & 0.996 & 0.997 & 0.998 \\
7 & 0.995 & 0.995 & 0.995 & 0.995 & 0.998 \\
8 & 0.779 & 0.785 & 0.737 & 0.814 & 0.814 \\
9 & 0.136 & 0.131 & 0.103 & 0.154 & 0.234 \\
10 & 0.495 & 0.532 & 0.400 & 0.663 & 0.721 \\
11 & 0.707 & 0.704 & 0.639 & 0.690 & 0.734 \\
12 & 0.768 & 0.757 & 0.746 & 0.774 & 0.790 \\
13 & 0.686 & 0.699 & 0.638 & 0.708 & 0.726 \\
14 & 0.947 & 0.953 & 0.963 & 0.955 & 0.962 \\
15 & 0.134 & 0.125 & 0.109 & 0.157 & 0.303 \\
16 & 0.353 & 0.392 & 0.238 & 0.552 & 0.632 \\
17 & 0.806 & 0.818 & 0.715 & 0.832 & 0.863 \\
18 & 0.811 & 0.815 & 0.776 & 0.831 & 0.838 \\
19 & 0.507 & 0.507 & 0.476 & 0.628 & 0.809 \\
20 & 0.653 & 0.659 & 0.743 & 0.654 & 0.698 \\
\bottomrule
\end{tabular}%
}
\label{tab:causal-tep-faults}
\end{table*}

\begin{table}[p]
\centering
\small
\caption{Detection and alarm outcomes at 50\% bandwidth with method-retrained detectors. TEP detection is binary (any fault), with event boundaries preserved within each run. Delay is conditional on detection; misses are reported separately. Entries are means and standard deviations across five joint experimental seeds.}
\resizebox{\textwidth}{!}{%
\begin{tabular}{@{}llccccccc@{}}
\toprule
Dataset & Method & Budget & Precision & Recall & FPR & Missed events & Delay (samples) & FA/1000 \\
\midrule
SMD-1-1 & PCA-Triage & 0.5 & 0.317 $\pm$ 0.013 & 0.507 $\pm$ 0.023 & 0.114 $\pm$ 0.003 & 2.4 $\pm$ 1.3 & 26.070 $\pm$ 17.873 & 6.159 $\pm$ 0.543 \\
SMD-1-1 & Variance & 0.5 & 0.329 $\pm$ 0.016 & 0.535 $\pm$ 0.033 & 0.114 $\pm$ 0.008 & 1.4 $\pm$ 1.1 & 14.523 $\pm$ 15.970 & 6.988 $\pm$ 0.456 \\
SMD-1-1 & Threshold & 0.5 & 0.320 $\pm$ 0.014 & 0.509 $\pm$ 0.017 & 0.113 $\pm$ 0.007 & 2.2 $\pm$ 1.3 & 28.300 $\pm$ 17.448 & 6.342 $\pm$ 1.036 \\
SMD-1-1 & Uniform & 0.5 & 0.319 $\pm$ 0.008 & 0.518 $\pm$ 0.009 & 0.116 $\pm$ 0.003 & 1.2 $\pm$ 1.1 & 10.901 $\pm$ 7.975 & 5.337 $\pm$ 0.234 \\
SMD-1-1 & Full Data & 1.0 & 0.330 $\pm$ 0.010 & 0.537 $\pm$ 0.011 & 0.114 $\pm$ 0.004 & 1.0 $\pm$ 1.0 & 4.967 $\pm$ 0.709 & 8.968 $\pm$ 0.552 \\
\addlinespace
PSM & PCA-Triage & 0.5 & 0.976 $\pm$ 0.031 & 0.016 $\pm$ 0.001 & 0.000 $\pm$ 0.000 & 62.0 $\pm$ 1.9 & 906.807 $\pm$ 359.620 & 0.046 $\pm$ 0.046 \\
PSM & Variance & 0.5 & 0.987 $\pm$ 0.014 & 0.016 $\pm$ 0.001 & 0.000 $\pm$ 0.000 & 61.0 $\pm$ 1.6 & 952.049 $\pm$ 214.155 & 0.036 $\pm$ 0.025 \\
PSM & Threshold & 0.5 & 0.980 $\pm$ 0.012 & 0.016 $\pm$ 0.001 & 0.000 $\pm$ 0.000 & 63.0 $\pm$ 2.0 & 985.411 $\pm$ 323.244 & 0.064 $\pm$ 0.042 \\
PSM & Uniform & 0.5 & 0.974 $\pm$ 0.022 & 0.017 $\pm$ 0.002 & 0.000 $\pm$ 0.000 & 61.4 $\pm$ 2.3 & 678.599 $\pm$ 402.486 & 0.068 $\pm$ 0.055 \\
PSM & Full Data & 1.0 & 0.963 $\pm$ 0.026 & 0.018 $\pm$ 0.004 & 0.000 $\pm$ 0.000 & 56.6 $\pm$ 1.1 & 629.185 $\pm$ 95.826 & 0.155 $\pm$ 0.096 \\
\addlinespace
TEP & PCA-Triage & 0.5 & 0.963 $\pm$ 0.001 & 0.904 $\pm$ 0.005 & 0.384 $\pm$ 0.012 & 0.0 $\pm$ 0.0 & 2.962 $\pm$ 0.879 & 7.223 $\pm$ 0.391 \\
TEP & Variance & 0.5 & 0.962 $\pm$ 0.001 & 0.903 $\pm$ 0.004 & 0.393 $\pm$ 0.009 & 0.0 $\pm$ 0.0 & 3.460 $\pm$ 0.554 & 7.170 $\pm$ 0.668 \\
TEP & Threshold & 0.5 & 0.965 $\pm$ 0.002 & 0.886 $\pm$ 0.007 & 0.354 $\pm$ 0.023 & 0.0 $\pm$ 0.0 & 4.588 $\pm$ 0.696 & 5.760 $\pm$ 0.398 \\
TEP & Uniform & 0.5 & 0.962 $\pm$ 0.001 & 0.914 $\pm$ 0.002 & 0.399 $\pm$ 0.011 & 0.0 $\pm$ 0.0 & 2.044 $\pm$ 0.438 & 9.672 $\pm$ 0.148 \\
TEP & Full Data & 1.0 & 0.943 $\pm$ 0.002 & 0.972 $\pm$ 0.002 & 0.638 $\pm$ 0.019 & 0.0 $\pm$ 0.0 & 0.350 $\pm$ 0.304 & 8.343 $\pm$ 1.114 \\
\addlinespace
\bottomrule
\end{tabular}%
}
\label{tab:detection}
\end{table}

At 50\% bandwidth on SMD, PCA-Triage obtains point F1 0.390; the highest constrained mean is 0.407 (Variance).
At 50\% bandwidth on PSM, PCA-Triage obtains point F1 0.031; the highest constrained mean is 0.033 (Uniform).
On all-20-fault TEP, PCA-Triage obtains weighted F1 0.675, versus 0.707 for Uniform and 0.728 for Full Data.
No inferential statistical superiority is claimed. Complete distortion, payload, delay, false-alarm and threshold-sensitivity tables are supplied in the accompanying supplementary material.

}{
The generated result table is unavailable. Run \texttt{make revised-tables}
before compiling the submission copy; legacy results must not be substituted.
}

\section{Discussion}

The PCA score should be interpreted as retained-subspace participation, not as
a direct measure of fault relevance or conditional redundancy. Consequently,
the method may be useful when dominant covariance structure aligns with
measurement utility, but variance or uniform allocation may be preferable when
low-variance channels contain rare diagnostic signals. Results are therefore
reported by dataset, fault/event, and bandwidth rather than collapsed into a
universal superiority statement.

The strict protocol updates PCA from reconstructed data and can accumulate
feedback bias. Periodic exploration, missing-data subspace tracking, or a local
full-rate gateway are possible extensions, but they change the information or
communication model and require separate evaluation. Similarly, delayed
linear interpolation can improve offline reconstruction but cannot support
zero-delay monitoring claims.

The component count has a precise interpretation in the revised formulation.
When $k=d$, the raw score is the marginal variance of the same accumulated
covariance. This identity does not make the smoothed historical method identical
to the preceding-window Variance baseline.
For $k<d$, it discounts variance lying outside the retained subspace. This may
remove nuisance variation, but it may also suppress an independent channel
whose low-variance behaviour is safety critical. A minimum target rate reduces complete
starvation in full windows but does not make the unsupervised score task-optimal. Conditional
redundancy or supervised relevance would require a different score and a
different claim.

Measurement and detection rankings need not agree. A strategy can reduce RMSE
by preserving high-energy signals yet remove a small diagnostic transient, or
increase RMSE while leaving a detector's decision boundary intact. For this
reason, the paper does not use F1 as a surrogate for measurement fidelity and
does not use reconstruction error as proof of fault-detection utility. Both are
reported alongside event recall and delay.

\section{Limitations}

The study uses replayed public data rather than a physical sensor network.
The low full-rate anomaly F1 scores limit conclusions about practical diagnostic
quality: the Isolation Forest is a fixed evaluation consumer, not an optimized
anomaly detector. Threshold sensitivity exposes this dependence but does not
establish deployment adequacy.
Payload counts omit protocol headers unless explicitly added for a deployment.
The public replay datasets provide no sensor-calibration uncertainty metadata,
so the study reports reconstruction distortion but not a metrological
uncertainty budget.
PCA captures second-order linear structure only. The score is unsupervised and
need not align with rare faults. Normal-only anomaly benchmarks require a
detector different from supervised TEP classification. Hardware timing and
energy measurements remain necessary before claiming embedded deployment
readiness.

Internal validity is improved by causal timing, common reconstruction, common
integer quotas, and train-only preprocessing, but results still depend on the
chosen downstream detectors and global parameter setting. Construct validity
is limited because payload values do not include packet headers, acknowledgments,
clock synchronization, or control messages. External validity is limited by
replayed benchmarks, one SMD machine, and the absence of a physical radio or
industrial fieldbus. These limitations motivate hardware validation and
multi-machine replication rather than stronger deployment language.

\section{Conclusion}

We formulated PCA-guided sensor-rate allocation as a causal measurement
problem with an explicit information pattern and hard per-window budget. The
revised method uses the diagonal of a rank-$k$ covariance approximation, makes
next-window decisions from past reconstructed observations, and avoids future-
value interpolation. Its value must be judged jointly by measurement
distortion, realized communication, event detection, delay, and false alarms.
The resulting formulation supports a narrower but testable contribution and
clearly identifies settings in which simpler allocation rules may be better.

\section*{Data and code availability}
The accompanying source archive contains the causal pipeline, official-partition
loaders, fixed seeds, regression tests, source/data SHA-256 manifests, measured
per-window acquisition ledgers, and scripts that regenerate all tables.
The accompanying archive defines the evaluated revision; repository details
are supplied on the separate author/title page. Public data are obtained
from their original repositories under their respective terms.

\clearpage
\bibliographystyle{unsrt}
\bibliography{references}
\end{document}
